Fix the leftmost gap of maximum height in each row \(P\). At that
gap order the crossing arcs from outermost to innermost. Use the
first \(m=r-1\) for the insertion, so that the choices for different
\(r\) are nested. Write \(K_{P,Q,r}\) for the value of the entry
after inserting \(t_{N,r}(XY)f^u_m\) and closing with \(Q\).
Put \(z=XY\) and
\[
                 E_r(z)=\prod_{\substack{d\mid r\\ d>2}} B_d(z).
\]

The complement of the box in the sphere is a disk. It contains
circles and a matching of the \(2m\) ports. Delete the circles starting
with any innermost disk on the side away from the box. Such a disk
contains no port or arc ending at a port. The deletion contributes
the colour weight of its interior face. The resulting monomial \(C\)
is common to all terms of the projector and to the identity box:
each term has the same prescribed boundary shading. After all these
deletions, an arc returning to the same row of ports has an innermost
adjacent turnback on that row, so the value with the projector is
zero. This uses the turnback relation over the fraction field;
the cleared box already has coefficients in the polynomial ring.

If every complementary arc passes from one row of ports to the other,
noncrossing forces the unique order-preserving matching. Choose
the exterior point next to the left edge of the box and draw this
matching as a right closure. Let \(w\in\{X,Y\}\) be the weight of the
left face, and \(w'\) the opposite weight, so that \(ww'=z\).
For the left boundary colour \(u\), Hazi's disk rule gives
\([2]_u=-w'\). Closing the rightmost strand keeps the left boundary
colour fixed. Iterating the partial trace identity
\(\operatorname{pTr}(f_j^u)=-[j+1]_u f_{j-1}^u/[j]_u\)
in the fraction field gives the planar scalar \((-1)^m[r]_u\)
\cite[equation (7)]{hazi}. Multiplying by \(w\) counts the outer
face once. The spherical values of the boxes alone are therefore
\[
 W_f=
 \begin{cases}
       z E_r(z),&r\ \text{even},\\
       w E_r(z),&r\ \text{odd},
 \end{cases}
 \qquad
 W_{\rm id}=w^{\lfloor m/2\rfloor+1}(w')^{\lceil m/2\rceil}.
\]
Indeed \([r]_u=E_r(z)\) for odd \(r\) and
\([r]_u=-w' E_r(z)\) for even \(r\). For the identity, the \(m\)
nested circles give alternating face colours starting with \(w\).
Choosing the exterior point on the other side gives the same two
spherical values: the side colours coincide for even \(m\), and both
expressions are symmetric in them for odd \(m\). The boundary ports
of the box stay fixed throughout. With the outside of the stack on
the left, \(w=Y\) for the chosen prefix: the outer face of the
matching meets the white interval from \(2n-1\) to \(0\).

Write \(a=|\pi\vee\rho|\) and \(a'=|\pi'\vee\rho'|\). Restoring
the identity box recovers the original pair, hence
\(C W_{\rm id}=J_n(P,Q)=X^aY^{a'}\). Combining the two values with
the definitions of \(t_{N,r}\) and \(b_k\) gives the exact quotient
\begin{equation}
\frac{K_{P,Q,r}}{b_k(X,Y)}
=
\begin{cases}
 X^{a-\lfloor r/2\rfloor}Y^{a'-\lfloor r/2\rfloor},
   &\text{the selected arcs lie in distinct cycles of }P\cup Q,\\
 0,&\text{otherwise}.
\end{cases} \label{eq:through-quotient}
\end{equation}
To see the criterion, cut the \(m\) selected arcs in the identity
diagram. When they belong to distinct cycles, the two ends of each
one are connected by the remainder of its own cycle. Conversely, the
unique through matching reconnects precisely those pairs, so
restoring the identity puts each selected arc on its own cycle.
In the through case the displayed quotient equals \(C\) for even
\(r\) and \(wC\) for odd \(r\); its exponents are nonnegative.
This proves the required integral divisibility for every pair.
